\documentclass[a4paper,11pt]{article}
\usepackage[utf8]{inputenc}
\usepackage[T1]{fontenc}
\usepackage[italian]{babel}
\usepackage{geometry}
\usepackage{graphicx}
\usepackage{booktabs}
\usepackage{amsmath}
\usepackage{hyperref}
\geometry{margin=2.5cm}

\title{ Pipeline di Preprocessing per Suoni Cardiaci PCG}
\author{Gianluca Farinaccio, Alexis Hidalgo Osvaldo Martinez}
\date{Febbraio 2026}

\begin{document}
\maketitle

\begin{abstract}
Questo lavoro presenta una pipeline completa per la classificazione binaria di suoni cardiaci fonocardiografici (PCG) in classi \emph{healthy} e \emph{unhealthy}. L'obiettivo principale non è solo addestrare una CNN, ma confrontare in modo sistematico diverse strategie di preprocessing audio prima della trasformazione in Mel-spectrogramma. Il progetto include: audit qualitativo e quantitativo del preprocessing, generazione del dataset preprocessato, training supervisionato con gestione dello sbilanciamento di classe e un modulo sperimentale di segmentazione cardiaca \emph{Springer-like} basato su HMM. I risultati preliminari mostrano accuratezza di validazione fino a $72.76\%$ con ResNet18, confermando la sensibilità del modello alle scelte di preprocessing.
\end{abstract}

\section{Introduzione}
L'analisi automatica dei suoni cardiaci può supportare lo screening precoce di anomalie valvolari o funzionali, in particolare in contesti a risorse limitate. In questo scenario, la classificazione PCG rappresenta un problema di machine learning con alcune criticità note: rumore ambientale, variabilità inter-soggetto, campionamenti eterogenei e sbilanciamento tra classi.

In questo progetto è stato utilizzato il dataset \emph{Heart Sound Database} pubblicato su Kaggle~\cite{heartsoundkaggle}. L'approccio adottato converte ciascun segnale audio in una rappresentazione tempo-frequenza e applica una rete neurale convoluzionale (ResNet18~\cite{he2016resnet}) per la classificazione binaria.

Il contributo principale del lavoro è il confronto tra pipeline di preprocessing diverse, con attenzione alla robustezza numerica (NaN/Inf, peak, RMS) e alla qualità della rappresentazione finale per il training.

\section{My Work}
Il progetto è stato sviluppato in Python con una struttura modulare composta da script eseguibili e libreria interna (`src/`), con configurazione delle pipeline in JSON.

\subsection*{Pipeline audio e rappresentazione}
Le pipeline sono definite in `configs/pipelines.json` come sequenze di step componibili (\texttt{load}, \texttt{bandpass}, \texttt{fft\_denoise}, \texttt{wavelet\_denoise}, \texttt{normalize\_peak/rms}, \texttt{trim/chunk}, \texttt{mel}, \texttt{log}).

I principali settaggi globali usati sono:
\begin{itemize}
  \item sample rate target: 2000 Hz;
  \item durata fissa: 10 s;
  \item Mel-spectrogramma: $n\_fft=256$, hop=128, 64 bande Mel, banda utile 20--300 Hz.
\end{itemize}

\subsection*{Audit del preprocessing}
Con lo script `scripts/audit_preprocessing.py` è stato eseguito un controllo su più pipeline:
\begin{itemize}
  \item validazione di shape e durata degli output;
  \item verifica assenza NaN/Inf;
  \item confronto di RMS e peak prima/dopo preprocessing;
  \item ispezione visiva waveform/FFT/Mel.
\end{itemize}

Dal file `runs/audit/summary.json` emerge che:
\begin{itemize}
  \item \texttt{bandpass\_mel}: \texttt{peak\_after\_max} = 0.99;
  \item \texttt{fft\_denoise\_mel}: \texttt{peak\_after\_max} = 0.967;
  \item \texttt{wavelet\_denoise\_mel}: \texttt{peak\_after\_max} = 0.980;
  \item \texttt{bandpass\_mel\_lowring}: \texttt{peak\_after\_max} = 1.413 (indicatore di overshoot/ringing residuo).
\end{itemize}

\subsection*{Preprocessing massivo e training CNN}
Con `scripts/data_preprocessing.py` sono stati preprocessati tutti i file train/val in PNG. La fase di training è stata implementata in `cnn/train.py` usando PyTorch~\cite{paszke2019pytorch}:
\begin{itemize}
  \item backbone ResNet18;
  \item input grayscale replicato su 3 canali, resize a 224x224;
  \item ottimizzatore Adam (default), 10 epoche;
  \item gestione sbilanciamento tramite \texttt{WeightedRandomSampler} e class weights quando rapporto classi $\geq 1.5$.
\end{itemize}

Nel run principale con dati preprocessati \texttt{bandpass\_mel}, il train set contiene 3240 campioni (2575 healthy, 665 unhealthy; rapporto $\approx 3.87$), mentre la validazione contiene 301 campioni.

\subsection*{Modulo di segmentazione cardiaca}
È stato inoltre sviluppato un modulo sperimentale in `src/segmentation.py` e `scripts/segmentation.py` che implementa una segmentazione \emph{Springer-like}:
\begin{itemize}
  \item envelope tramite trasformata di Hilbert + low-pass;
  \item stima picchi ed BPM;
  \item HMM a 4 stati (S1, sistole, S2, diastole);
  \item decoding Viterbi in log-domain;
  \item visualizzazione segmenti su segnale.
\end{itemize}
Questo modulo è utile come base per integrare informazioni di fase cardiaca in futuri modelli.

\section{Risultati}
La Tabella~\ref{tab:cnn_results} riporta i risultati disponibili nei run di training salvati.

\begin{table}[h]
\centering
\caption{Risultati di validazione (runs/cnn).}
\label{tab:cnn_results}
\begin{tabular}{lccc}
\toprule
Run & Pretrained & Best epoch & Best val acc \\
\midrule
2026-01-24 11:17 & True  & 6  & 0.7276 \\
2026-01-28 15:04 & False & 10 & 0.7176 \\
\bottomrule
\end{tabular}
\end{table}

Osservazioni preliminari:
\begin{itemize}
  \item il modello raggiunge accuratezza di validazione massima del 72.76\%;
  \item l'uso di pesi preaddestrati migliora leggermente il picco di accuratezza, ma con instabilità nelle ultime epoche;
  \item la qualità del preprocessing influenza in modo significativo la distribuzione energetica dello spettrogramma e quindi la classificazione.
\end{itemize}

\section{Sviluppi Futuri}
\begin{itemize}
  \item confronto sistematico tra tutte le pipeline (\texttt{bandpass}, \texttt{fft\_denoise}, \texttt{wavelet}) a parità di seed e iperparametri;
  \item introduzione di metriche più robuste allo sbilanciamento (F1 macro, recall classe \emph{unhealthy}, AUC);
  \item aggiunta di calibration e threshold tuning per ridurre falsi negativi clinicamente critici;
  \item integrazione della segmentazione HMM nel preprocessing (feature per fase cardiaca);
  \item valutazione di architetture alternative (EfficientNet, CRNN, Transformer audio);
  \item validazione cross-dataset per testare generalizzazione e robustezza.
\end{itemize}

\section{Bibliografia}
\bibliographystyle{plain}
\bibliography{references}

\end{document}
